\documentclass[10pt,a4paper]{article}

\usepackage[T1]{fontenc}
\usepackage[utf8]{inputenc}
\usepackage[ngerman]{babel}
\usepackage[margin=1.5cm]{geometry}
\usepackage{lmodern}
\usepackage{tabularx}
\usepackage{enumitem}
\usepackage[hidelinks]{hyperref}

\pagestyle{empty}
\setlength{\parindent}{0pt}
\setlength{\parskip}{0pt}
\setlist[itemize]{leftmargin=1.2em, nosep, topsep=2pt}
\urlstyle{same}

% Einheitliche Gestaltung; Abstände hier zentral anpassen.
\newcommand{\cvsection}[1]{%
  \par\addvspace{10pt}%
  {\large\bfseries #1}\par\nobreak
  \vspace{2pt}\hrule height 0.5pt\relax
  \vspace{5pt}%
}

\newcommand{\cvemployer}[2]{%
  \noindent\begin{tabularx}{\linewidth}{@{}>{\raggedright\arraybackslash}X r@{}}
    \textbf{#1} & #2 \\ 
  \end{tabularx}\par\nobreak
}

\newcommand{\cvrole}[2]{%
  \noindent\begin{tabularx}{\linewidth}{@{}>{\raggedright\arraybackslash}X >{\raggedleft\arraybackslash}p{4cm}@{}}
    \textit{#1} & \textit{#2} \\ 
  \end{tabularx}\par\nobreak
}

\newcommand{\cventry}[4]{%
  \cvemployer{#1}{#2}%
  \cvrole{#3}{#4}%
}

\begin{document}

\begin{center}
  {\LARGE\bfseries Jordi Buscaglia}\par
  \vspace{5pt}
  {\small Am Elisabethgehölz 6, 20535 Hamburg\par
  \vspace{2pt}
  +49\,157\,53646577\enspace\textbar\enspace
  \href{mailto:Jordi.buscaglia577@gmail.com}{Jordi.buscaglia577@gmail.com}}
\end{center}

\cvsection{Berufserfahrung}

% Bitte prüfen: Beginn der Masterarbeit; Datum aus der Vorlage übernommen.
\cventry{Fraunhofer-Institut für Siliziumtechnologie}{Itzehoe}
  {Masterand im Bereich Mikroelektronische Systeme}{Jun. 2026 -- heute}
\begin{itemize}
  \item Forschung zu Retrieval-Augmented Generation (RAG) mit Fokus auf die Verbesserung des Tabellen-Parsings und die Implementierung neuer Verfahren.
  \item Machbarkeitsstudien zur Codegenerierung mit lokalen Sprachmodellen wie Qwen, Llama und Mistral.
\end{itemize}
\vspace{0.5cm}

\cvrole{Werkstudent im Bereich optische Systeme}{Sep. 2025 -- Jun. 2026}
\begin{itemize}
  \item Forschung zu Retrieval-Augmented Generation (RAG) mit Fokus auf die Verbesserung des Tabellen-Parsings und die Implementierung neuer Verfahren.
  \item Machbarkeitsstudien zur Codegenerierung mit lokalen Sprachmodellen wie Qwen, Llama und Mistral.
\end{itemize}
\vspace{0.5cm}


% Bitte prüfen: Zeitraum der Bachelorarbeit; Datum aus der Vorlage übernommen.
\cvemployer{Philips Healthcare}{Hamburg}
\cvrole{Bachelorand im Bereich Advanced Research and Development}{Jan. 2025 -- Sep. 2025}
\begin{itemize}
  \item Bachelorarbeit: \glqq Large Language Models für die Softwareentwicklung: Evaluierung der Machbarkeit lokaler Large Language Models zur Codegenerierung\grqq.
\end{itemize}
\vspace{0.5cm}

\cvrole{Praktikant im Bereich Advanced Research and Development}{Okt. 2024 -- Jan. 2025}
\begin{itemize}
  \item Forschung zu Retrieval-Augmented Generation (RAG) mit Fokus auf die Verbesserung des Tabellen-Parsings und die Implementierung neuer Verfahren.
  \item Machbarkeitsstudien zur Codegenerierung mit lokalen Sprachmodellen wie Qwen, Llama und Mistral.
\end{itemize}
\vspace{0.5cm}

\cvrole{Werkstudent in der Testautomatisierung von Röntgendetektoren}{Jan. 2024 -- Sep. 2024}
\begin{itemize}
  \item Erstellung und Durchführung von Testplänen und Dauerlauftests; Anpassung der Teststände, Protokollierung von Messdaten und Nachbearbeitung von Röntgenaufnahmen zur Qualitätssicherung.
  \item Erfahrung im technischen Projektmanagement und in der Zusammenarbeit im internationalen Konzernumfeld.
\end{itemize}

\cvsection{Bildungsweg}

\cventry{Hochschule für Angewandte Wissenschaften}{Hamburg}
  {Mikroelektronische Systeme (M.Sc.)}{2025 -- heute}
\vspace{5pt}

\cventry{Hochschule für Angewandte Wissenschaften}{Hamburg}
  {Mechatronik (B.Sc.), Schwerpunkt Elektrische Antriebe und Robotik}{2020 -- 2025}
\vspace{5pt}

\cventry{Margaretha-Rothe-Gymnasium}{Hamburg}
  {Leistungskurse: Mathematik, Physik, Biologie und Englisch}{2011 -- 2019}

\cvsection{Kenntnisse}

\begin{tabularx}{\linewidth}{@{}l >{\raggedright\arraybackslash}X@{}}
  \textbf{Programmierung:} & Python, C \\[3pt]
  \textbf{Software:} & MATLAB, Simulink, Altium, LTspice \\[3pt]
  \textbf{Weitere Werkzeuge:} & Git, GitHub, Codex
\end{tabularx}

\cvsection{Sprachen}

\textbf{Deutsch:} Muttersprache\quad
\textbf{Englisch:} Verhandlungssicher\quad
\textbf{Spanisch:} Grundkenntnisse

\end{document}
